\documentclass[11pt]{article}
\usepackage[a4paper,margin=0.95in]{geometry}
\usepackage{amsmath,amssymb,amsthm,hyperref}
\newtheorem{theorem}{Theorem}
\newtheorem{lemma}{Lemma}
\title{A failure of the DeTemple complete-monotonicity threshold}
\author{Research proof and Lean formalization}
\date{October 27, 2026}
\begin{document}
\maketitle
\begin{abstract}
We refute the sufficient condition in Yang's complete-monotonicity conjecture,
retained in the published paper of Yang, Chu and Zhang. The failure occurs for
a positive parameter strictly inside the proposed range. The proof uses local
continuity of the digamma function and a logarithmic midpoint obstruction.
The literal source refutation, including the interior-parameter conclusion, is
formalized in Lean.
\end{abstract}

For $x>0$, let
\[
 \psi(x)=\frac{\Gamma'(x)}{\Gamma(x)},\qquad
 R(x)=\psi(x+\tfrac12)-\log x,
\]
and use the exact function from \cite[eq.~(3.3)]{Yang}
\[
 G_a(x)=24\left(x^4+ax^2+\frac7{40}a-\frac{31}{336}\right)R(x)
          -\left(x^2-\frac7{40}+a\right).
\]
Conjecture~1(ii) of that paper states that $-G_a$ is completely monotone on
$(0,\infty)$ if and only if $a\le155/294$. Complete monotonicity means smoothness
and $(-1)^nf^{(n)}(x)\ge0$ for every $n\ge0$ and $x>0$; in particular it implies
convexity.

\begin{theorem}\label{counter}
There exist $0<a<155/294$ and $x>0$ such that
\[
 G_a(2x)<\frac{G_a(x)+G_a(3x)}2.
\]
Consequently $-G_a$ is not convex on $(0,\infty)$ and is not completely monotone.
Thus the conjectured sufficient condition is false.
\end{theorem}

The strict inequalities for both $a$ and $x$ are part of the conclusion. The ratio
monotonicity question in Conjecture~1(i) is not settled here.

\begin{lemma}\label{local}
Let $c>0$ and let $h$ be real-valued and continuous at zero. Put
\[
 F_c(x)=24(x^4+cx^2)(h(x)-\log x)-x^2,\quad x>0.
\]
For all sufficiently small positive $x$,
$F_c(3x)-2F_c(2x)+F_c(x)>0$.
\end{lemma}
\begin{proof}
For $x>0$, expanding and using $\log(kx)=\log k+\log x$ gives
\begin{equation}\label{difference}
 \frac{F_c(3x)-2F_c(2x)+F_c(x)}{x^2}
 =24\bigl(B_c(x)-(50x^2+2c)\log x\bigr)-2,
\end{equation}
where
\begin{align*}
 B_c(x)={}&(81x^2+9c)h(3x)-2(16x^2+4c)h(2x)+(x^2+c)h(x)\\
          &-(81x^2+9c)\log3+2(16x^2+4c)\log2.
\end{align*}
Continuity gives
$B_c(x)\to2ch(0)-9c\log3+8c\log2$ as $x\downarrow0$.
Since $50x^2+2c\to2c>0$ and $-\log x\to+\infty$, the right side of
\eqref{difference} tends to $+\infty$. Multiplying by $x^2>0$ proves the claim.
\end{proof}

\begin{proof}[Proof of Theorem~\ref{counter}]
Set $a_0=155/294$. The constant term in the polynomial multiplying $R$ cancels:
\[
 \frac7{40}a_0-\frac{31}{336}=0.
\]
The function $h(x)=\psi(x+1/2)$ is continuous at zero because $\Gamma$ is
analytic and nonzero near $1/2$. Therefore
\[
 G_{a_0}(x)=F_{a_0}(x)+\frac7{40}-a_0
\]
satisfies the strict midpoint inequality for some $x>0$ by Lemma~\ref{local}.
For this fixed $x$, the midpoint defect
\[
 a\longmapsto\frac{G_a(x)+G_a(3x)}2-G_a(2x)
\]
is continuous (indeed affine) in $a$ and is positive at $a_0$. It remains positive
at some $a$ with $0<a<a_0$. Negating the inequality contradicts convexity of
$-G_a$, proving both conclusions.
\end{proof}

\section*{Formal scope}
The formal definition of $\psi$ is the actual real derivative
\texttt{deriv Real.Gamma / Real.Gamma}. Its local continuity is proved from
complex Gamma analyticity, with the real derivative identification checked.
The proof includes \eqref{difference}, its one-sided infinite limit, the
strictly positive midpoint witnesses, continuity in the source parameter and
the implication from complete monotonicity to convexity. The endpoint
\texttt{DeTemple.}\allowbreak\texttt{source\_counterexample\_328867ii} is a premise-free negation
of the literal source claim; \texttt{source\_interior\_not\_}\allowbreak\texttt{completelyMonotone}
also records $0<a<155/294$.

The certificate uses Lean~4.34.1 and mathlib at the pinned revision
\begin{center}
\texttt{d13f23b723b8a846827a245b89c10fc7d3f11612}.
\end{center}
The separate rational point in the working mathematical notes is not asserted
as a formal endpoint. This manuscript claims no optimal replacement threshold.
Standard Gamma, convexity, continuity and logarithm facts are credited
ingredients. No worldwide-priority claim is made.

\begin{thebibliography}{9}
\bibitem{Yang} Z.-H. Yang, Y.-M. Chu and X.-H. Zhang,
\emph{Necessary and sufficient conditions for functions involving the psi
function to be completely monotonic}, Journal of Inequalities and Applications
\textbf{2015}, 157; \href{https://doi.org/10.1186/s13660-015-0674-8}{doi:10.1186/s13660-015-0674-8}.
The conjecture also appears in Z.-H. Yang, arXiv:1409.6411 (2014),
\emph{Several completely monotone functions related to DeTemple's sequence}.
\end{thebibliography}
\end{document}
